\section{Redes Neuronales Artificiales}
Una red neuronal artificial es un modelo de computación basado en las redes neuronales biológicas.

\subsection{Neurona Artificial}
Una red neuronal está compuesta por neuronas, que de manera similar a sus análogos biológicos forman
conexiones con otras neuronas transmitiendo activaciones a través de dichas conexiones. 
Una neurona realiza dos funciones fundamentales:

\begin{description}
    \item[Función de propagación] Recibe los valores de salida de otras neuronas y los combina 
    considerando el peso de sus conexiones para obtener la entrada neta de la neurona.
    Una de las funciones más utilizadas es la suma ponderada, en la que se suman los productos 
    del valor de salida de cada neurona de entrada y el peso de su conexión, dada por:
    \begin{equation}
        \text{net} = \sum_{i=1}^{n}\left(x_{i}\cdot w_{i}\right)
    \end{equation}
    donde
    \begin{description}
        \item[$x_{i}$] es el valor de salida de la $i$-ésima neurona de entrada.
        \item[$w_{i}$] es el peso asociado a la conexión de la $i$-ésima neurona de entrada.
        \item[$n$] es el número total de neuronas de entrada.
    \end{description}
    \item[Función de activación] Transforma la entrada neta en un valor de activación.
\end{description}
%% Mencionar bias

Estas dos funciones permiten modelar el comportamiento básico de una neurona artificial dentro de una red neuronal.
%% Agregar diagrama de una neurona

\subsection{Funciones de Activación}
La elección de la función de activación en una red neuronal es un aspecto fundamental, 
ya que influye directamente en la capacidad del modelo para representar y aproximar funciones complejas. En particular, estas funciones introducen no linealidad en la red, lo cual es esencial para modelar relaciones no lineales entre los datos de entrada.

Entre las funciones de activación más utilizadas se encuentran las funciones sigmoides y 
las funciones lineales por partes.

% Mencionar linearidad y no-linearidad
% Ejemplos: Sigmoide, ReLU, tanh


\subsection{Arquitectura de Redes Neuronales}

% Capas: entrada, oculta, salida
% Feedforward flow
% Diagrama

\subsection{Proceso de aprendizaje}
% Funcion de perdida
% Backpropagation
% Descenso de gradiente


\subsection{Limitaciones de Redes Totalmente Conectadas}
% Muchos parametros
% Falta de espacialidad
% Escalabilidad pobre para imagenes

\section{Redes Neuronales Convolucionales}
\section{Técnicas de detección de objetos}
\section{YOLOv8}
\section{Interpretabilidad}
\section{Visualización}
\section{Blender}